
\documentclass{article}
\pdfpagewidth=8.5in
\pdfpageheight=11in
\usepackage{ijcai20}

% Use the postscript times font!
\usepackage{times}
\renewcommand*\ttdefault{txtt}
\usepackage{soul}
\usepackage{url}
\usepackage[hidelinks]{hyperref}
\usepackage[utf8]{inputenc}
\usepackage[small]{caption}
\usepackage{graphicx}
\usepackage{amsmath}
\usepackage{booktabs}
\urlstyle{same}

\usepackage{todonotes}

\begin{document}

\title{
    \includegraphics[width=8cm, keepaspectratio]{tudelftlogo.png}\\
    \vspace*{2cm}
    \textbf{
        \textless Title of your paper\textgreater\\
        {\large \textless Subtitle of your paper\textgreater}
    }\\
    \vspace*{1cm}
}

\author{
    \textless \textbf{
Your Name$^1$\textgreater}\\
    \hfill \break
    \textbf{Supervisor(s): \textless Responsible Professor$^1$\textgreater, \textless Supervisor$^1$\textgreater }\\
    \break
%    \affiliations
    {\large 
        \hfill \break
        $^1$EEMCS, Delft University of Technology, The Netherlands
    }\\
}

\date{}

\maketitle
\thispagestyle{empty}

\let\clearpagebackup\clearpage
\renewcommand{\clearpage}{ }

\onecolumn

\vspace*{1.5cm}
\begin{center}
    A Thesis Submitted to EEMCS Faculty Delft University of Technology,\\
    In Partial Fulfilment of the Requirements\\
    For the Bachelor of Computer Science and Engineering\\
    \today
\end{center}

\vspace*{2cm}

\noindent
{\small
Name of the student: \textless Your name\textgreater\\
Final project course: CSE3000 Research Project\\
Thesis committee: \textless Responsible Professor\textgreater, \textless Supervisor\textgreater, \textless Examiner\textgreater\\
}
\vfill

\begin{center}
    An electronic version of this thesis is available at http://repository.tudelft.nl/.
\end{center}

\twocolumn
\let\clearpage\clearpagebackup  
\clearpage
\setcounter{page}{1}


\begin{abstract}
Although transformers are state-of-the-art models for natural language tasks, obtaining reasonable performance still often requires large transformers which are expensive to train and deploy. Fortunately, there are techniques to increase the size of transformers without extra computing costs. One such technique is sparsity. However, it remains unclear whether sparse architecture is intrinsically more efficient than its dense counterpart. In this paper, we investigate whether replacing the feedforward networks in small transformers with sparse alternatives results in better predictions and faster inference. We found that although inference speed does not increase due to software and hardware limitations, certain sparse alternatives do result in better language understanding. Our research contributes to smarter architectural decision making when designing small language models.

% The aim of this template is to make it more clear what is expected from you.
% \textbf{It is by no means required to follow this exact same structure.}
% The abstract should be short and give the overall idea:
% what is the background, the research questions, what are your contributions, and what are the main conclusions.
% It should be readable as a stand-alone text (preferably no references to the paper or to outside literature).
\end{abstract}

\section{Introduction}
\label{sec:intro}
\begin{itemize}
    \item Sparse transformers are transformer models that only activate a portion of their parameters when processing an input.
    \item Sparse transformers such as GLaM, Switch Transformer, and Scaling Transformer have been shown to outperform dense models such as GPT-3 \cite{du_glam_2022}, T5 \cite{fedus_switch_2022}, and BIGBIRD-RoBERTa \cite{jaszczur_sparse_2021}.
    \item But the architectural reason for the sparse transformer's success is uncertain.
    \item Sparsity allows the model more parameters without extra computational costs. More parameters bring greater learning capacities.
    \item But when such models are compared to much smaller baselines, it becomes unclear whether the improved performance is inherently related to sparsity, in the sense of more efficient utilization of model parameters, or simply the result of a bigger model.
    \item Works on network pruning demonstrated that it is indeed possible to achieve better \cite{balderas_optimizing_2024} or at least comparable performance \cite{lecun_optimal_1989} to the original model with a sparsified network.
    \item More theoretical backings of sparse feedforward networks \cite{csordas_approximating_2023,baykal_theoretical_2022}
    \item Our work contributions to the question of sparsity's influence in learning capacity by comparing sparse transformers against dense transformers of similar sizes.
    \item \textbf{TODO: findings and implications}

\end{itemize}

The remainder of this paper is organized as follows: In section \ref{sec:background} we provide an overview of sparse transformers. Section \ref{sec:approach} elaborates on the sparse feedforward paradigms by deriving their sizes. These results are used in section \ref{sec:setup}, which describes the experiment setup. Our findings are presented in section \ref{sec:results}. This is followed by an analysis and interpretability study in section \ref{sec:discussion}. In section \ref{sec:responsible} we reflect on the ethics and reproducibility of our research. Lastly, we state the conclusion and prescribe future directions for our research in section \ref{sec:conclusion}.



\section{Methodology, Background, Problem Description}
Choose one that fits your research best:
\subsection{Methodology and/or background}
Typically in general research articles, the second section contains a description of the research methodology, explaining what you, the researcher, is doing to answer the research question(s), and why you have chosen this method.
For purely analytical work this is a description of the data collection or experimental setup on how to test the hypothesis, with a motivation.

In any case this section includes references to necessary background information.
For a survey paper this includes the method of how you arrived at the set of papers included in the survey.

\subsection{Formal Problem Description}
For some types of work in computer science the methodology is standard: analyze the problem (e.g., make assumptions and derive properties), present a new algorithm and its theoretical background, proving its correctness, and evaluate unproven aspects in simulation.
Then an explanation of the methodology is often omitted, and the setup of the evaluation is part of a later section on the evaluation of the ideas.\footnote{This already shows that there is no single outline to be given for all papers.}
In this case, explain relevant (background) concepts, theory and models in this section (with references) and relate them to your research question.
Also this section then typically contains a more precise, formal description of the problem.

Do not forget to give this section another name, for example after the problem you are solving.


\section{Approach}
\label{sec:approach}
\begin{itemize}
    \item Assume $d_\text{intermed} = Md_\text{in}$
    \item Let $d$ denote dimension and $N$ denote quantity
    \item For vanilla FFN, total number of parameters is $2d_\text{intermid}d_\text{in}$
    \item For MoE, total number of parameters is $2d_\text{in}N_\text{experts}(d_\text{intermed}+1)$
    \item For Controller FFN, total number of parameters is $2d_\text{in}(d_\text{intermed}+\frac{d_\text{intermed} + d_\text{in}}{2N_\text{blocks}})$
    \item For product key memory, total number of parameters is $2d_\text{in}(N_\text{head}d_\text{subkey}+\frac{N_\text{subkey}^2}{2}+\frac{N_\text{head}N_\text{subkey}d_\text{subkey}}{d_\text{in}})$
    \item $d_\text{in}$ is fixed, the other hyperparameters are set so that sparse FFNs has the same number of parameters as vanilla FFN.
\end{itemize}




\section{Experimental Setup and Results}
As discussed earlier, in many sciences the methodology is explained in section 2 and this section only discusses the results. 
However, in computer science, most often the details of the evaluation setup are described here first (simulation environment, etc.).
Very important is that any skilled reader would be able to reproduce this setup and then obtain the same results.

Then, results are reported in an accessible manner through figures (preferably with captions that allow them to be understood without going through the whole text), observations are made that clearly follow from the presented results.
Conclusions are drawn that follow logically from the previous material.
Sometimes the conclusions are in fact hypotheses, which in turn may give rise to new experiments to be validated.

You may want to give this section another name.

\section{Responsible Research}
\label{sec:responsible}
Deep learning, among other fields, faces a reproducibility crisis \cite{semmelrock_reproducibility_2023}. We combated this issue in our research by providing a replication package containing all the source code and instructions to reproduce the results presented in this paper. Furthermore, we reported our hyperparameters and hardware in \Cref{sec:setup}. All our models were pretrained and evaluated with fixed seeds. These techniques together ensure that our experiments and findings can be replicated exactly.

In accordance with the Netherlands Code of Conduct for Research Integrity, we reported all our experiment results in this paper under the principles of honesty and transparency. In addition, following the standards for good research practices (chapter 3), we refrained from data fabrication and manipulation. For example, our pretraining and evaluation sets are separate and all models were given to the same data.

Our work raises ethical concerns as it is related to language models. Language models may be trained on copyrighted or private texts. In addition, when trained with domain knowledge, language models can be misused by malicious actors. To mitigate these issues, all our models were pretrained on a synthetic dataset. The dataset contains no knowledge of any specific domain.

% Reflect on the ethical aspects of your research and discuss the reproducibility of your methods.
% Note that although in many published works there is no such a section (it may be part of some meta-information collected by the journal, or part of the discussion section), we require you to think (and report) about this as part of this course.

\section{Discussion}
\label{sec:discussion}
\subsection{Performance of Sparse Feedforward Layers}
Whether sparse feedforward layers can achieve comparable performance as the vanilla FFN depends on the baseline architecture and type of sparse feedforward layer as \Cref{fig:spar-rat,fig:spar-rat-super} demonstrate. Regardless, sparse models seem to be fast learners since the loss difference between the sparse models and the baseline is much smaller by the end of pretraining than at the start. This is likely because sparse feedforward layers approximate a larger network than an FFN of the same size. Sparse feedforward layers are more flexible learners. In terms of neural memories, although MoE, CNT, and PKM do not have as many unique value vectors as FFN, their memory coefficients are computed from more complex, non-linear functions.

% all types of sparse feedforward layers except PKM for BERT models, outperformed the RoBERTa baseline on BLiMP tasks. This is not the case for GPT-Neo-based models though. CNT and PKM can achieve comparable performance to the baseline in certain configurations, but overall sparse feedforward layers led to worse performance. In should be noted that in general, all models in this experiment performed poorly on BLiMP tasks. None were significantly better than a random classifier.
% \Cref{fig:eval-loss} presents a similar finding. Although no sparse models achieved lower evaluation loss than their respective baselines, among the GPT-Neo-based models, the MoE variant achieved similar evaluation loss towards the end of pertaining. At any rate, sparse models seem to be fast learners, as the loss difference between the sparse models and the baseline is much smaller by the end of pertaining than at the start. This is corroborated by the RoBERTa-based models. Despite having a higher loss at initialization, almost all sparse models reached a similar evaluation loss as the baseline after the first epoch.

% CNT is consistently the best type of sparse feedforward layer for BLiMP tasks. This is possibly due to CNT approximating a much larger network than the other types of sparse feedforward layers in the experimented configurations. If CNT, MoE, and PKM are all interpreted as mixture of expert models, then, thanks to weight sharing, the CNT models in this experiment contain combinatorially more experts than the MoE and PKM models. It would also imply that for a mixture of expert models, the number of experts is more important to performance than the number of activated experts and the intermediate size of each expert; for CNT, only the top expert is activated and the intermediate size is much smaller than in the input size.

\subsection{Role of Sparsity Ratio}
There is inconclusive evidence on the number of active parameters directly affecting performance. Even when considering a sparse variant such as MoE in isolation, it is not guaranteed for performance to decrease or increase with the sparsity ratio. A plausible explanation is that the feedforward architecture is far more important to language understanding than the proportion of active parameters. In other words, how parameters are activated is far more crucial than the amount of active parameters. This hypothesis is corroborated by the CNT results. CNT outperformed the baselines on both BLiMP and (Super)GLUE tasks despite having less than $5\%$ parameter activation in the feedforward layer. In contrast, PKM cannot compete with the BLiMP baselines even when more than $60\%$ of its parameters are active because it is not as architecturally efficient as CNT.
 
% number learnable parameters is more important to performance than the number of active parameters. The models investigated in this paper are large enough such that they can learn the data distribution with only a fraction of their parameters. This hypothesis is corrobrated by the CNT results. CNT outperformed the baselines on both BLiMP and (Super)GLUE tasks despite having less than $5\%$ parameter activation in the feedforward layer.

% On BLiMP tasks, there is a moderate negative correlation between the sparsity and performance for RoBERTa-based models, as sparser models tends to score higher. The relationship is reversed, albeit with weaker correlation, for GPT-Neo based models, where denser models seems to perform better. The architecture is likely to play a more important role. In other words, how parameters are activated is far more crucial to performance than the amount of active parameters. For example, using CNT, sparse transformers can reach similar BLiMP scores as the baseline, despite having less than $5\%$ parameter activation in the feedforward layer, whereas with PKM, sparse transformers cannot compete against the baseline even when more than $60\%$ of the feedforward parameters are active. This result suggests that CNT is a more efficient architecture than PKM.

\subsection{Inference Cost of Sparse Models}
% The vanilla feedforward network has by far shorter inference time than its sparse variants as illustrated by \Cref{tab:inference}. 
Despite the theoretical speedups sparsity and conditional computation bring, sparse feedforward layers are slowed down by software limitations and the small transformer setting. All models are implemented with PyTorch \cite{paszke_pytorch_2019}, which, at time the of writing, lacks support for various essential features on sparse tensors, such as reshapes and automatic differentiation, resulting in inefficient implementations. Furthermore, the speedups of conditional computation often manifest in a distributed setting, as is the case in \cite{fedus_switch_2022}, where an MoE model is scaled up by distributing the experts across multiple devices. Small transformers do not enjoy this advantage because all training happens on one GPU.

\subsection{Threats to Validity}
Threats to the validity of this work can be categorized into three categories: threats to internal validity, which are the confounders of the experiment, threats to external validity, which are the issues surrounding the generalizability of the experiment, and threats to construct validity, which concerns the soundness of the research questions.
\subsubsection{Internal Validity}
While the sizes of the sparse feedforward layers are set to be as close as possible to their dense counterparts in the experiments, it is not possible for them to be identical. Nonetheless, the relative difference is insignificant so it is unlikely to result in a huge change in learning capacity. Other hyperparameters may have a significant influence on the experiment result. For example, we hypothesize based on \Cref{fig:eval-loss} that all sparse models will converge to the same evaluation loss as that of the baselines given more epochs. However, this would require more computing budget than what is available for this research.
\subsubsection{External Validity}
It must be emphasized that the model configurations investigated in the experiments are by no means exhaustive. As showcased by \Cref{tab:mode-size}, despite considering different configurations for each type of sparse feedforward layer, the sparsity ratio is in some cases almost identical. This is because for some models, such as PKM, a noticeable change in sparsity ratio requires a huge increase in parameter count. 

All experiments are conducted on a single GPU. While this is unlikely to affect pretraining and evaluation, changing the hardware setup will alter the inference speed of sparse transformers, especially if the models are distributed over multiple devices.
\subsubsection{Construct Validity}
There is no universal measure of sparsity. In this research, our measure of sparsity differs per model. For MoE, the degree of sparsity is defined as the number of experts selected. For, CNT, the degree of sparsity is measured in the number of blocks the activation vector is split into. For PKM, the degree of sparsity is defined as the number of keys selected. These definitions do not always correspond to the sparsity ratio well. Although the sparsity ratio appears to be a universal definition of sparsity, it too has shortcomings. As an example, PKM's sparsity ratio is challenging to measure as it depends on the input. Even under the neural memory framework, the number of active memory cells is not guaranteed to be static, so it too is a difficult sparsity measure to implement.

\subsection{Future Work}
\Cref{tab:summary_eval} shows that CNT is an efficient architecture. It outperforms the baselines on both BLiMP and (Super)GLUE while using the lowest percentage of parameters. This paper investigated CNT under a very low sparsity ratio. For future research, we recommend studying the performance of CNT with more active parameters. For other models, we recommend reducing the intermediate size, and instead increasing the number of experts for MoE or keys for PKM. We conjecture that a very small intermediate size does not necessarily lead to worse performance, as shown by CNT under the mixture of experts interpretation. Lastly, we recommend adding more layers and controllers to CNT to mimic hierarchical MoE. Weight sharing the MoE and PKM is also a good direction for further research.

\section{Conclusions}
\label{sec:conclusion}
This paper investigates whether sparse feedforward layers in small transformers result in better language understanding and faster inference than the standard feedforward network. We found that inference speed does not increase due to software and hardware limitations associated with our small language model setting. Some sparse feedforward layers, such as the controller feedforward layer, in the right configurations, achieves better language understanding than the standard feedforward network while activating just a fraction of their parameters. We hypothesize that this is because the sparse feedforward layers more efficiently utilize their parameters by approximating larger networks. We recommend future research exploring more configurations of different sparse feedforward architectures.


% Briefly summarize the (main) research question(s).
% Provide your conclusions, the answers to the research question(s).
% Make statements.
% Highlight interesting elements, contributions.

% Discuss open issues, possible improvements, and new questions that arise from this work; formulate recommendations for further research.

% Ideally, this section can stand on its own: it should be readable without having read the earlier sections and accessible to anyone with a bachelor degree in Computer Science.


\appendix
\section{Some further guidelines that go without saying (right?)}

\begin{itemize}
\item Read the manual for the Research Project. (See e.g.\ the instructions on the maximum length: less is more!)
\end{itemize}

\subsection{Reference use}
\begin{itemize}
\item use a system for generating the bibliographic information automatically from your database, e.g., use BibTex and/or Mendeley, EndNote, Papers, or \ldots
\item all ideas, fragments, figures and data that have been quoted from other work have correct references
\item literal quotations have quotation marks and page numbers
\item paraphrases are not too close to the original
\item the references and bibliography meet the requirements
\item every reference in the text corresponds to an item in the bibliography and vice versa
\end{itemize}

\subsection{Structure}
Paragraphs
\begin{itemize}
\item are well-constructed
\item are not too long: each paragraph discusses one topic
\item start with clear topic sentences
\item are divided into a clear paragraph structure
\item there is a clear line of argumentation from research question to conclusions
\item scientific literature is reviewed critically
\end{itemize}

\subsection{Style}
\begin{itemize}
\item correct use of English: understandable, no spelling errors, acceptable grammar, no lexical mistakes 
\item the style used is objective
\item clarity: sentences are not too complicated (not too long), there is no ambiguity
\item attractiveness: sentence length is varied, active voice and passive voice are mixed
\end{itemize}

\subsection{Tables and figures}
\begin{itemize}
\item all have a number and a caption
\item all are referred to at least once in the text
\item if copied, they contain a reference
\item can be interpreted on their own (e.g. by means of a legend)
\end{itemize}


\bibliographystyle{plain}
\bibliography{references}

A rule of thumb for dealing with the literature is the following: scan about 10--20 contributions: read title, abstract, part of introduction and conclusions; categorize contribution; some of these are studied in more depth: completely read about 5 conference papers or equivalent (summarize contribution in own words); of which studied in-depth about 2 conference papers (the student is able to explain in detail and criticize contributions). This may result in 5--20 references, possibly even more if the project is a literature study.


\end{document}

